\section{8. Engineering Data
Repository}\label{engineering-data-repository}

\subsection{8.1 Identification and
Scope}\label{identification-and-scope}

The GSEG engineering data repository contains controlled definitions of
the segment design, particularly parameters, configurations, source
documents and links among them. \textbf{It is not the operational DAS
archive}, which stores EWAS mission data.

Project-documentation sources are maintained in the GitHub repository
\texttt{papakallo/\allowbreak{}ewas}, including \texttt{GSEG/\allowbreak{}TRS/\allowbreak{}} and
\texttt{GSEG/\allowbreak{}DDF/\allowbreak{}}. Version control of source files does not in itself
constitute a complete register of all engineering parameters. GSEG
therefore provides for \textbf{one controlled project parameter
register} linked to document versions.

\subsection{8.2 Engineering Data
Collections}\label{engineering-data-collections}

\begingroup\GSEGTableStyle{3}
\begin{longtable}[]{|P{0.28}|P{0.40}|P{0.32}|}
\hline\rowcolor{TableHeaderBlue}
\begin{minipage}[b]{\linewidth}\raggedright
\color{white}\bfseries\sffamily
Data collection
\end{minipage} & \begin{minipage}[b]{\linewidth}\raggedright
\color{white}\bfseries\sffamily
Location / reference
\end{minipage} & \begin{minipage}[b]{\linewidth}\raggedright
\color{white}\bfseries\sffamily
Status
\end{minipage} \\
\hline
\endhead
\hline
\endlastfoot
Segment requirements and identifiers & \texttt{GSEG/\allowbreak{}TRS/\allowbreak{}}, Chapter 7 &
Exists; seven-subsystem structure covered by TRS \\ \hline
Functional and physical architecture & \texttt{GSEG/\allowbreak{}DDF/\allowbreak{}}, Chapters 4--5
& Preliminary DDF definition \\ \hline
Technical allocations and margins & \texttt{GSEG/\allowbreak{}DDF/\allowbreak{}}, Chapter 6 &
Preliminary design allocations, including \textbf{100 ppm GSEG
unavailability}: 90 ppm allocated and 10 ppm reserve \\ \hline
Central GSEG parameter register &
\path{GSEG/DDF/engineering_data/GSEG_parameters_register.csv} &
\textbf{Preliminary register present in repository; not approved} \\ \hline
GSEG--SSEG and inter-subsystem interface definitions & Relevant
ICDs/IRDs & Agreement and controlled issues required \\ \hline
Subsystem design data and justifications & Future subsystem TRS/DDF/DDJF
documents and GSEG DDJF & Not yet complete \\*
\end{longtable}
\endgroup

\newpage
\subsection{8.3 Parameter Identification and
Status}\label{parameter-identification-and-status}

Each GSEG parameter-register entry should identify, at a minimum:
\textbf{stable parameter identifier, name, responsible element, value
with unit, status, source requirement or decision, reference version and
confirmation method}.

The \textbf{B} and \textbf{P} statuses are defined in Section 1.3.
\textbf{A/V --- Assessed/Verified} denotes a value supported by
controlled analysis or verification results; its mere entry in the
register does not confer this status.

Links between parent parameters and allocations to GSN, C\&C, DPS, EWS,
MCC, DAS and GIS must permit assessment of the impact of a value change
on the design and interfaces. Supporting allocation data remain in the
DDJF; the register provides their identifier or location rather than
repeating calculations.

\subsection{8.4 Completeness Status}\label{completeness-status}

The GSEG engineering data repository is \textbf{not declared complete}
at this stage. Detailed ICDs, physical-item configuration identifiers,
lower-level allocations and verification evidence in particular remain
to be brought under control.

The register accompanying this working version constitutes an initial
list of selected segment parameters, not an approved parameter baseline
for all seven subsystems. It will be extended and agreed as lower-level
documentation develops.
